% ============================================================
% Part III §III.2 — The Boredom Experiment Re-Read — DRAFT v1 (FCDP v2)
% Pack: plans/packs/part-iii-2-pack-v1.md (D1 embedded) · Quotes: part-iii-2-quotes.json
% M1 = Band R (the record); M2–M6 = Band G. Footnote-wrapped cites at write time.
% ============================================================

\section{The Boredom Experiment Re-Read}

The record comes first, plainly stated, with its reliability flags attached. The first study was a within-participants physiological pilot with three participants. Each watched three seventeen-minute videos: a boring control (a 1989 word-processor tutorial), an interesting control (an alien-conspiracy video), and the clinical stimulus (360-degree footage of spinal-deformation surgery). An electroencephalographic cap recorded four active channels (F3, F4, P3, P4) at 200 Hz, targeting the default-mode network's alpha (8--12 Hz) and theta (4--8 Hz) bands; heart rate, gaze, and pupil diameter came from the headset's sensor suite. Its clean findings are two. All three participants felt the clinical video as longer than its seventeen minutes---the one clean recording's owner felt it as twenty-five to thirty---and self-reported engagement ran high, with two participants rating interest at eight of eight. Its headline physiological finding is suggestive only: the clinical video's cortical profile more closely resembled the boring control than the interesting one, but only one of the three recordings was clean, one carrying movement artifacts and one impedance problems. The heart measures were inconclusive, and learning itself, despite the study's title, was never measured. The second study replaced the brain with the eye. Twelve participants watched the same three videos while gaze was tracked at 120 Hz; the analysis took pupil dilation, gaze position, deviation from center, and---its strongest measure---the variance of gaze deviation under a five-point sliding-window median, with analysis of variance at $\alpha = 0.05$. A four-person webcam sub-study validated the approach against consumer hardware (average cross-correlation 0.874). The presentation order was changed to interesting, then clinical, then boring, because a boring-first order had left participants too exhausted for what followed---a fact that will matter later. The findings: gaze-variance on the clinical video was statistically distinct from the boring video ($p = 0.0004$) and resembled the interesting one, while every other feature-by-video comparison stayed nonsignificant; pupil dilation exceeded the boring baseline for the interesting video in six of twelve participants; gaze-variance ran lower for interesting than boring in eight of twelve; and the study drew its own distinction between two boredoms---a high-arousal boredom whose gaze roves the screen searching for something to attend to, and a low-arousal boredom, the zombie stare, in which variance collapses, observed in two participants. Around the published studies stands the raw archive of the same program, reported here with its unpublished status marked and its participants anonymized. Its instrument is a seven-item questionnaire answered after each video: fatigue before; boredom and engagement, each on a nine-point scale; minutes until boredom began; a sleep-fight rating; felt duration; fatigue after. Eight participants completed it for all three videos. The aggregates: mean boredom of 7.0 for the boring video, 6.0 for the clinical, 4.5 for the interesting; felt duration of twenty minutes or more, against seventeen actual, in five of eight for the boring video, six of eight for the clinical, five of eight for the interesting---sixteen of twenty-four episodes dilated, and every felt value in this record that cleared the clock also cleared twenty minutes; boredom onset within five minutes in seven of eight for the boring video, two of them under forty seconds; and fatigue after the boring video at or above its starting level in six of eight.

Read through the frame, the boring stimulus behaves exactly as becoming-bored-by predicts, and its behavior calibrates everything else. The first form of boredom has a nameable culprit---this tutorial, this drone---and a proto-definition that the lecture course states with both structural moments in one breath: ``that which bores, which is boring, is that which holds us in limbo and yet leaves us empty.''\footnote{Martin Heidegger, \textit{The Fundamental Concepts of Metaphysics: World, Finitude, Solitude}, trans. William McNeill and Nicholas Walker (Bloomington: Indiana University Press, 1995), 87 / GA 29/30: 130--32.} Both moments have empirical faces here. Being left empty (\textit{Leergelassenheit}) shows as the collapse of interest ratings toward the floor and the near-immediacy of onset: seven of eight participants were bored within five minutes of a seventeen-minute film, two of them within forty seconds, which is to say that the tutorial's offer was refused almost at the threshold, and everything after the refusal was endurance. Being held in limbo (\textit{Hingehaltenheit}) shows in the sleep-fight column and in the stretched while---twenty and twenty-five and thirty felt minutes against the seventeen actual, one participant reaching a full felt hour. And between the two moments runs the visible counter-move that the eye-tracker caught as high variance: the searching gaze, roving the screen for anything to attend to, is \textit{Zeitvertreib} made measurable, the passing-of-time in its restless first-form costume. The lecture course had already named the mechanism whereby the counter-move fails: ``Passing the time is a driving away of boredom that drives time on.''\footnote{Heidegger, \textit{Fundamental Concepts}, 94 / GA 29/30: 141--42.} Passing the time drives time on and thereby keeps the dragging while in view; the search sustains the very emptiness it flees; and the eye that cannot find anything is, in its motion, the portrait of an offer that gives nothing to complete.

The clinical stimulus is another matter, and it is the section's center, because on it the instruments disagree. The cortical profile of the first study read the surgical footage as boring-like; the gaze of the second study read it as engaging-like with the corpus's one emphatic significance; the participants called themselves interested; and the clean recording's owner felt the seventeen minutes as twenty-five to thirty. The frame reads the disagreement as the signature of the second form, being-bored-\textit{with}, whose whole structure is an occupied surface over an emptied depth [a projected reading, and marked as such]. In the lecture course's paradigm---an agreeable dinner party after which one finds, at home, that one was bored after all---the passing-of-time no longer sits beside the activity as cigarette or doodle; the occupation itself becomes the passing-of-time: ``It is not smoking as an isolated occupation, but our entire comportment and behaviour''\footnote{Heidegger, \textit{Fundamental Concepts}, 112 / GA 29/30: 169--70.} An engaged-looking gaze over a hollow cortical hour is precisely that structure rendered in channels: the surface busy, the depth left empty, and no culprit anywhere to blame, for the second form's emptiness, as the course insists, ``what is boring does not come from outside: it arises from out of Dasein itself.''\footnote{Heidegger, \textit{Fundamental Concepts}, 128 / GA 29/30: 192--94.} The objection writes itself, and it deserves its strength: the two instruments belong to two studies, two years, two cohorts, and their disagreement might be an artifact of the seam. The concession is real and it is bounded, twice over. First, the keystone does not need the electroencephalogram: within the later study's own logic, felt duration against gaze supplies the two channels, and a gaze that reads engaged over a while that stretches past its clock is the same divergence without the seam. Second---and this is the raw archive's singular contribution---one participant staged the divergence inside a single instrument. Asked two questions instead of one, the participant coded S07 in the archive rated the clinical film at boredom six \textit{and} engagement seven, and the interesting film at boredom seven \textit{and} engagement six. Given separate voices for surface and depth, the answer sheet refused the bipolar axis on its own initiative; the two layers, permitted to speak separately, spoke differently. That is the two-channel rule confirmed from inside testimony itself, and it is why no instrument built on a single engaged-to-bored dimension could ever have seen what this record shows.

The felt-duration column deserves its own reading, because it is the corpus's cleanest depth evidence and its most Heideggerian. Sixteen of twenty-four episodes were felt as longer than their seventeen minutes, and the dilation crosses every stimulus category: the boring film stretches, the clinical film stretches, and the interesting film stretches for five of eight participants too. Boredom's temporality, the course argues, is not a matter of the clock running slowly but of the while itself lengthening for the one held in it---``In boredom, the while [\textit{Weile}] becomes long [\textit{lang}].''\footnote{Heidegger, \textit{Fundamental Concepts}, 152 / GA 29/30: 229.}---and the drag is a relation to time rather than a property of the timed: the interval oppresses us ``not because the progress of time is slow, but because it is \textit{too} slow.''\footnote{Heidegger, \textit{Fundamental Concepts}, 97 / GA 29/30: 146--47.} What the column shows, read so, is that the laboratory situation as such---strapped into a headset, held before a mandated film with nothing to do but watch---is already a holding apparatus, a small architecture of the held-in-limbo whose grip tightens or loosens by stimulus but never quite releases. The admission must then be paid for, because it presses on the keystone's fallback: if dilation alone were the depth criterion, the interesting film---engaged-reading gaze in eight of twelve, stretched for five of eight---would qualify as well. The fallback therefore stands on a conjunction, not on the stretch alone: an engaged-reading surface over a depth that both lengthens \textit{and} rates itself bored, which the clinical film satisfies at a mean of six against the interesting film's four and a half---and, at its sharpest, on the answer sheet that carried both poles at once. The extremity makes the structure vivid: one participant, bored within twenty-five seconds, felt the seventeen minutes as an hour, and for the same participant the clinical film did the opposite work, restoring fatigue from six to three---the while releasing where interest took hold. What is boring, at the limit, is not this tutorial or that surgery. The course's terminal thesis names it: ``What is boring is neither beings nor things as such [\ldots] but temporality as such''\footnote{Heidegger, \textit{Fundamental Concepts}, 158 / GA 29/30: 236.}

At the record's low-arousal pole stands the zombie stare, and with it the analysis must grow careful, because the third form of boredom enters this Part only as an analogy and is marked as one. In two participants the searching variance collapsed: the gaze stopped roving, the counter-move ceased, and the posture resembled---resembled---the stance the course reserves for profound boredom, in which all passing-of-time has become powerless and what remains is ``a being compelled to listen''\footnote{Heidegger, \textit{Fundamental Concepts}, 136 / GA 29/30: 205.} No claim follows that a participant underwent a metaphysical disclosure between films; the category caution forbids it, and the data do not need it. What the record adds instead is the exit that the analogy alone would miss. One participant fell asleep during the boring film, rated the sleep-fight at nine, and felt the seventeen minutes as five---the column's most extreme contraction. Where dilation is the signature of wakeful holding, sleep is not boredom's maximum but its escape: the entranced while requires a witness who remains held, and when the witness slips away, the long while vanishes with the sleeper. The methodological consequence is concrete, and it is double. A felt-duration probe saturates strangely at the sleep boundary, so any classifier built on these signatures needs sleep-exit as its own terminal category, distinct from the deepest boredom rather than continuous with it; and a collapsed gaze is by itself ambiguous between absorption and abandonment, so no classifier may read low variance alone---the companion channels, sleep-fight and felt time, must disambiguate what the stilled eye cannot.

Two features of the record remain, and they close the laboratory reading by turning it against the schema that most naturally organizes such experiments. The first is carryover. The second study reordered its films because boring-first had exhausted its participants, and the raw archive confirms the pattern from inside: fatigue after the boring film met or exceeded its starting level in six of eight participants, twice climbing from a three to a seven and an eight. Boredom, whatever else it is, depletes---an aftermath the lecture course, which treats the attunement as disclosive, never theorizes, and the one place in this corpus where the traffic runs from the laboratory back toward the philosophy. Design will have to respect it: pacing, sequence, and washout are boredom variables, not logistics. The second feature is the anomaly pair, and it is the record's quiet refutation of stimulus-essentialism. One participant entered the clinical film already exhausted, at fatigue seven, and rated it boredom nine and engagement one, the film felt as thirty minutes; the same participant then met the boring tutorial---the corpus's designated dullness---and rated it boredom one, engagement seven, felt eight minutes, fatigue falling from seven to one. Another participant was bored by the interesting film within a minute while giving the tutorial an engagement of five. The same footage bores, engages, restores, and repels, depending on the episode's temporal and dispositional situation; nothing in the stimulus fixes its effect. The cause-effect schema fails empirically exactly where the lecture course dismantles it conceptually, and the failure is the reanalysis's warrant: what the instruments were watching was never a property of three films, but the fortunes of a chain---solicited, held, emptied, completed, or abandoned---whose structure the next sections carry out of the laboratory and into a deployed world.

% ============================================================
% FOOTER (FCDP provenance)
% Material: ALL NEW (greenfield; P7 = N/A). Retained/corrections: n/a. Omissions: none.
% Quotes: Q1–Q8 all used, footnote-wrapped at write time.
%   NOTE — drafting initially mis-numbered six markers against the bank; caught
%   pre-substitution and corrected (bank IDs authoritative): M3 uses Q4+Q5;
%   M4 uses Q7, Q3, Q8; M5 uses Q6. Lesson for next packs: draft with the bank
%   open, verify IDs before substitution.
% Counterargument dispositions: C1 ANSWER (M1 flags + M6 structural claims);
%   C2 CONCEDE-AND-LIMIT (M3 ¶: seam conceded; EEG-independent fallback + S07);
%   C3 ANSWER (M6 inversions run against expectation); C4 surfaced (M1 status marked).
% GAUNTLET REPORT (final): G-A PASS (avg 29.9, short .222, long .444; §III.0
%   waivers apply) · G-B PASS (8/8, exit 0, zero residue; ID-drift self-caught
%   pre-substitution, 6 markers renumbered) · G-C PASS (8 footnoted dual loci;
%   no memory loci) · G-D PASS (greps clean; "within-participants" fix;
%   Leergelassenheit/Hingehaltenheit first-use glosses added) · G-E PASS after
%   1 cycle (pack v1.1: E6 extended w/ S06 + fatigue climbs, E6a boring felt-
%   duration column; prose: "most extreme contraction"; dinner-party phrase
%   de-quoted to indirect speech) · G-F attested · G-G PASS after 1 cycle
%   (KEY FIX: M4 now pays for the cross-category-dilation admission — the
%   fallback restated as a CONJUNCTION [stretch ∧ elevated boredom rating,
%   6.0 vs 4.5] + the within-instrument case; felt-15 removed from stretch
%   evidence; 25–30 de-pluralized here AND in §III.0; zombie-ambiguity
%   disambiguation added to M5). R-cycles used: 2 of 3 (incl. pre-subst fix).
%   UNDER-TARGET NOTE: ~2,200w vs plan's 2,900–3,400 — all movements present;
%   expansion candidates at review: M2 (searching-gaze phenomenology), M5.
% ============================================================
